\clearpage
\section*{Problem 8}
\addcontentsline{toc}{section}{Problem 8}
\begingroup
\hypersetup{colorlinks=true,
            linkcolor=blue!50!black,
            citecolor=blue!50!black,
            urlcolor=blue!50!black}
\newcommand{\qeightR}{\mathbb{R}}
\newcommand{\qeightZ}{\mathbb{Z}}
\newcommand{\qeightSm}{\mathscr S}
\newcommand{\qeightLdual}{L^{\vee}}
\newcommand{\qeightdd}{\mathrm d}
\newcommand{\qeightLag}{\Lambda}
\newcommand{\qeightHess}{\operatorname{Hess}}
\newtheorem{qeighttheorem}{Theorem}
\newtheorem{qeightproposition}{Proposition}
\newtheorem{qeightlemma}{Lemma}
\newtheorem{qeightcorollary}{Corollary}
\theoremstyle{definition}
\newtheorem{qeightdefinition}{Definition}
\newtheorem{qeightremark}{Remark}
\title{Smoothing Polyhedral Lagrangian Surfaces}
\author{}
\date{}
\maketitle


\section*{Statement and Answer}

Equip $\qeightR^4$ with coordinates $(q_1,q_2,p_1,p_2)$ and the standard symplectic
form $\omega=\qeightdd p_1\wedge \qeightdd q_1+\qeightdd p_2\wedge \qeightdd q_2$. A \emph{polyhedral
Lagrangian surface} $K\subset\qeightR^4$ is a finite polyhedral complex, all of whose
faces are Lagrangian, that is a topological submanifold of $\qeightR^4$.

\medskip
\noindent\textbf{Theorem.} \emph{If $K$ is a polyhedral Lagrangian surface in
which exactly $4$ faces meet at every vertex, then $K$ admits a Lagrangian
smoothing: a Hamiltonian isotopy $K_t$ of smooth Lagrangian submanifolds
parametrized by $t\in(0,1]$, extending to a topological isotopy on $[0,1]$ with
$K_0=K$.}

\medskip
\noindent The answer is \textbf{YES}.

\medskip
The proof has a local part (a normal form near each vertex and edge, from which
a smooth structure on $K$ is read off) and a global part (a fibration of
Lagrangian planes transverse to $K$, with respect to which all sufficiently
$C^1$-small ``smoothing functions'' yield genuine graphical Lagrangians). The
$4$-faces hypothesis enters in the local normal form (Lemma~\ref{lem:linear}),
where four Lagrangian quadrants force a standard symplectic basis; it is not a
Maslov/parity condition.

\section*{1. Local model at a vertex}

\begin{qeightlemma}\label{lem:linear}
Let $\qeightR^2\to\qeightR^4$ be an embedding that is linear on each of the four quadrants,
with Lagrangian image $\Sigma$, and assume $\Sigma$ is not contained in a single
plane. Then there is a linear symplectic transformation of $\qeightR^4$ carrying
$\Sigma$ to the product of the union of the positive coordinate axes in
$\qeightR^2_{p_1q_1}$ with that in $\qeightR^2_{p_2q_2}$.
\end{qeightlemma}

\begin{proof}
Let $(v_1,v_2,u_1,u_2)$ be the tangent vectors at the origin to the four edges of
$\Sigma$, ordered so that cyclically adjacent vectors span the faces. The
pairings $\omega(v_i,u_i)$ cannot vanish: if some $\omega(v_i,u_i)=0$, then
$\omega$ would vanish identically on a $3$-dimensional subspace of $\qeightR^4$, which
is impossible for a symplectic form. After swapping $v_i\leftrightarrow u_i$ if
needed we may take both pairings positive, and after rescaling we may take
$\omega(v_1,u_1)=\omega(v_2,u_2)=1$ with the cross pairings vanishing (the faces
are Lagrangian). Then $(v_1,v_2,u_1,u_2)$ is a standard symplectic basis, and the
linear map $\partial_{p_i}\mapsto v_i$, $\partial_{q_i}\mapsto u_i$ is the
required symplectomorphism, carrying $\Sigma$ to the stated product of
positive-axis unions.
\end{proof}

In the plane $\qeightR^2_{pq}$ the symplectic pairing projects the union of the two
positive axes homeomorphically onto the dual of the antidiagonal line $p=q$.
Taking the product and applying Lemma~\ref{lem:linear}:

\begin{qeightcorollary}\label{cor:Ldual}
There is a linear Lagrangian plane $L\subset\qeightR^4$ such that the symplectic
pairing $\qeightR^4\to \qeightLdual$ restricts to a homeomorphism $\Sigma\to \qeightLdual$.
\end{qeightcorollary}

This equips $\Sigma$ with a smooth structure (pulled back from $\qeightLdual$), which
we fix henceforth. Given $L$, call a Lagrangian $\qeightLag\subset\qeightR^4$
\emph{graphical} if the pairing restricts to a diffeomorphism $\qeightLag\cong\qeightLdual$.

\section*{2. Smoothing functions}

Choose a Lagrangian splitting $\qeightR^4\cong L\oplus\qeightLdual$. Each of the four
quadrant-faces of $\Sigma$ is graphical over $\qeightLdual$, hence the graph of the
differential of a quadratic form; collecting these gives quadratic forms
$\{q_{ij}\}_{i,j\in\pm}$ on $\qeightLdual$, agreeing to first order along the edges.
Via $\Sigma\cong\qeightLdual$ (Corollary~\ref{cor:Ldual}) write $q_\Sigma$ for the
resulting $C^1$ function on $\Sigma$ that is $q_{ij}$ on each face.

\begin{qeightdefinition}\label{def:smoothing}
The space $\qeightSm(\Sigma)$ of \emph{smoothing functions} is the set of $C^1$
functions $f:\Sigma\to\qeightR$ such that $f+q_\Sigma$ is $C^\infty$.
\end{qeightdefinition}

$\qeightSm(\Sigma)$ is invariant under adding smooth functions, and depends only on
$L$, not on the splitting (a different splitting changes $q_{ij}$ by a global
quadratic form, i.e.\ a smooth function). To a smoothing function $f$ we
associate the Lagrangian $\qeightLag_{\qeightdd f}\subset\qeightR^4$ that is, piecewise, the graph
of $\qeightdd f$ over each face — equivalently the graph of $\qeightdd(f+q_\Sigma)$ in the
splitting.

\begin{qeightlemma}\label{lem:bijection}
$f\mapsto \qeightLag_{\qeightdd f}$ is a bijection between graphical Lagrangians and
smoothing functions on $\Sigma$ modulo additive constants.
\end{qeightlemma}

\begin{proof}
In the splitting, $\qeightLag_{\qeightdd f}$ is the graph of $\qeightdd(f+q_\Sigma)$ as a function
on $\qeightLdual$; graphical Lagrangians over $\qeightLdual$ are exactly graphs of
differentials of smooth functions, and $f+q_\Sigma$ is smooth precisely when
$f\in\qeightSm(\Sigma)$. The construction does not use the splitting, so the bijection
depends only on $L$.
\end{proof}

\noindent The same statements hold verbatim for an edge: a pair of Lagrangian
half-planes meeting along a line, with the analogue of
Corollary~\ref{cor:Ldual} provided by the next lemma.

\begin{qeightlemma}\label{lem:edge}
If $\Sigma$ is a pair of linear Lagrangian half-planes meeting along a line
$\ell$, the space of Lagrangian planes $L$ for which the pairing
$\Sigma\to\qeightLdual$ is a homeomorphism is contractible.
\end{qeightlemma}

\begin{proof}
Up to affine linear symplectomorphism, $\Sigma$ is the product of the real axis
in one $\qeightR^2$-factor with the union of the positive axes in another. The pairing
$\Sigma\to\qeightLdual$ is a homeomorphism iff $L$ is transverse to both half-planes,
i.e.\ the symplectic reduction of $L$ along $\ell$ is a line in $\ell^\perp/\ell\cong\qeightR^2$
meeting the interior of the positive quadrant — a contractible condition. The
space of Lagrangian lifts of such a line is itself contractible (two lifts differ
by the graph of a map $\ell'\to\ell$), giving the claim.
\end{proof}

\section*{3. Existence of small smoothing functions}

The crux is that $\qeightSm(\Sigma)$ is \emph{not} invariant under rescaling, so the
existence of small smoothing functions is not automatic. We obtain them by
\emph{mollifying} the $C^1$ reference function $q_\Sigma$. The estimate that makes
this quantitative is the following standard fact; we isolate it because the rate
$O(\varepsilon)$ in $C^1$ (not merely $o(1)$) is what the assembly in Section~6
needs, and it relies on $q_\Sigma$ having a \emph{Lipschitz} gradient, which is
true here for a structural reason.

\begin{qeightlemma}[Mollification of a function with Lipschitz gradient]\label{lem:moll}
Let $u\in C^1(\qeightR^n)$ have a globally $\Lambda$-Lipschitz gradient, i.e.
$|\nabla u(x)-\nabla u(y)|\le\Lambda|x-y|$ for all $x,y$. Fix an even, nonnegative
$\psi\in C^\infty_c(\qeightR^n)$ with $\int\psi=1$, and set
$\psi_\varepsilon(x)=\varepsilon^{-n}\psi(x/\varepsilon)$ and
$u_\varepsilon:=u*\psi_\varepsilon$. Then $u_\varepsilon\in C^\infty(\qeightR^n)$,
$(\varepsilon,x)\mapsto u_\varepsilon(x)$ is jointly smooth on $(0,\infty)\times\qeightR^n$,
and, with $M_1:=\int|s|\,\psi(s)\,\qeightdd s$ and $M_2:=\int|s|^2\psi(s)\,\qeightdd s$,
\[
\|\nabla u_\varepsilon-\nabla u\|_{C^0}\ \le\ \Lambda M_1\,\varepsilon,
\qquad
\|u_\varepsilon-u\|_{C^0}\ \le\ \tfrac12\Lambda M_2\,\varepsilon^2 .
\]
\end{qeightlemma}

\begin{proof}
Smoothness of $u_\varepsilon$ and joint smoothness in $(\varepsilon,x)$ are the
usual properties of convolution against a smooth compactly supported kernel that
depends smoothly on $\varepsilon>0$. Since $\nabla u_\varepsilon=(\nabla u)*\psi_\varepsilon$
and $\int\psi_\varepsilon=1$,
\[
\nabla u_\varepsilon(x)-\nabla u(x)
=\int\bigl(\nabla u(x-s)-\nabla u(x)\bigr)\,\psi_\varepsilon(s)\,\qeightdd s,
\]
whence, by the Lipschitz bound and the substitution $s=\varepsilon\sigma$,
\[
|\nabla u_\varepsilon(x)-\nabla u(x)|
\le\Lambda\!\int|s|\,\psi_\varepsilon(s)\,\qeightdd s
=\Lambda\varepsilon\!\int|\sigma|\,\psi(\sigma)\,\qeightdd\sigma=\Lambda M_1\varepsilon .
\]
For the function bound, write $u(x-s)-u(x)=-\nabla u(x)\!\cdot\! s+R(x,s)$ with
$R(x,s)=\int_0^1\bigl(\nabla u(x)-\nabla u(x-\theta s)\bigr)\!\cdot\! s\,\qeightdd\theta$,
so $|R(x,s)|\le\tfrac12\Lambda|s|^2$ by the Lipschitz bound. As $\psi$ is even,
$\int(\nabla u(x)\!\cdot\! s)\psi_\varepsilon(s)\,\qeightdd s=0$, and
\[
|u_\varepsilon(x)-u(x)|=\Bigl|\int R(x,s)\psi_\varepsilon(s)\,\qeightdd s\Bigr|
\le\tfrac12\Lambda\!\int|s|^2\psi_\varepsilon(s)\,\qeightdd s=\tfrac12\Lambda M_2\varepsilon^2 .
\qedhere
\]
\end{proof}

\begin{qeightlemma}\label{lem:small}
There is a smooth family $\{f^\varepsilon\}_{\varepsilon\in(0,1]}\subset\qeightSm(\Sigma)$
with $\|f^\varepsilon\|_{C^1}=O(\varepsilon)$. More precisely
$g^\varepsilon:=f^\varepsilon+q_\Sigma$ is smooth, $(\varepsilon,x)\mapsto
g^\varepsilon(x)$ is jointly smooth on $(0,1]\times\Sigma$, and there is a constant
$M$ with $\|\qeightdd f^\varepsilon\|_{C^0}\le M\varepsilon$ for all $\varepsilon\in(0,1]$.
In particular $\qeightSm(\Sigma)$ contains functions of arbitrarily small $C^1$-norm.
\end{qeightlemma}

\begin{proof}
Recall from Section~2 that, in the linear coordinates on $\Sigma\cong\qeightLdual\cong\qeightR^2$,
the reference function $q_\Sigma$ is \emph{piecewise quadratic}: it equals the
quadratic form $q_{ij}$ on the $(i,j)$-th face, and the four pieces agree to first
order along the shared edges (this first-order agreement is exactly the assertion
that $q_\Sigma$ is $C^1$, Section~2). Consequently the gradient $\nabla q_\Sigma$ is
\emph{continuous and piecewise linear}: on each face it is the linear map
$x\mapsto\qeightHess(q_{ij})\,x$, and these linear pieces match continuously across the
edges. A continuous, piecewise-linear vector field with finitely many linear pieces
is globally Lipschitz; let $\Lambda:=\mathrm{Lip}(\nabla q_\Sigma)<\infty$ (one may
take $\Lambda=\max_{ij}\|\qeightHess(q_{ij})\|$, since for $x,y$ in possibly different
faces one interpolates along the segment $[x,y]$, on which $\nabla q_\Sigma$ is
piecewise linear with slopes bounded by this maximum, and the continuity across
edges lets the increments telescope).

Thus $q_\Sigma$ has a globally $\Lambda$-Lipschitz gradient, and
Lemma~\ref{lem:moll} applies (with $n=2$). Fix an even bump $\psi$ as there and set
\[
g^\varepsilon:=q_\Sigma*\psi_\varepsilon\in C^\infty(\Sigma),
\qquad
f^\varepsilon:=g^\varepsilon-q_\Sigma .
\]
Then $f^\varepsilon+q_\Sigma=g^\varepsilon$ is $C^\infty$, so
$f^\varepsilon\in\qeightSm(\Sigma)$ by Definition~\ref{def:smoothing}. The map
$(\varepsilon,x)\mapsto g^\varepsilon(x)$ is jointly smooth on $(0,1]\times\Sigma$
by Lemma~\ref{lem:moll}. Finally Lemma~\ref{lem:moll} gives
$\|\qeightdd f^\varepsilon\|_{C^0}=\|\nabla g^\varepsilon-\nabla q_\Sigma\|_{C^0}
\le\Lambda M_1\varepsilon$ and
$\|f^\varepsilon\|_{C^0}\le\tfrac12\Lambda M_2\varepsilon^2$, so
$\|f^\varepsilon\|_{C^1}\le M\varepsilon$ with $M:=\Lambda(M_1+\tfrac12 M_2)$
(using $\varepsilon\le1$). This proves the quantitative claims, and letting
$\varepsilon\to0$ gives smoothing functions of arbitrarily small $C^1$-norm.
\end{proof}

\begin{qeightremark}
The previous version of this lemma built $f^\varepsilon$ by gluing
plane-models with a dilated partition of unity and asserted, without proof, that
the gluing region is where the face data agree to order $\varepsilon^2$. That
$O(\varepsilon^2)$ claim would require \emph{second-order} agreement of the face
quadratics along edges, whereas only \emph{first-order} ($C^1$) agreement holds;
the cutoff gradients $O(1/\varepsilon)$ would then defeat the bound. Mollification
sidesteps this entirely: it never enlarges gradients, and the $O(\varepsilon)$ rate
comes solely from the gradient of $q_\Sigma$ being Lipschitz, which the $C^1$,
piecewise-quadratic structure guarantees outright.
\end{qeightremark}

\section*{4. Global construction: a conormal fibration}

To globalize, we replace the single plane $L$ by a smoothly varying family.

\begin{qeightdefinition}\label{def:fibration}
A \emph{conormal fibration dual to $K$} is a smoothly varying family $L_z$ of
affine Lagrangian planes, $z\in K$, such that: (i) near each vertex the $L_z$ are
translates of a single plane as in Corollary~\ref{cor:Ldual}; (ii) $L_z$ varies
smoothly along edges and along faces; and (iii) in a collar of each edge the
planes in the normal direction are translates of those along the edge, with
\emph{matched} inward normal fields on the two adjacent faces (opposite vectors
under the identification $T_z\sigma\cong\qeightLdual_z\cong T_z\sigma'$).
\end{qeightdefinition}

\begin{qeightlemma}\label{lem:fibexists}
$K$ admits a conormal fibration $\{L_z\}_{z\in K}$ which near each vertex agrees
with the local choice of Corollary~\ref{cor:Ldual}. Writing $\ell_z:=L_z-z$ for the
linear part of the affine plane $L_z$ (so $L_z=z+\ell_z$ and $z\in L_z$), the family
can moreover be taken with the following \emph{local normal forms}, which we record
for the disjointness argument.
\begin{enumerate}[label=\textup{(N\arabic*)}]
\item\label{N:vertex} \emph{(Vertices.)} There is $r_0>0$ so that on
$B(v,r_0)\cap K$ about each vertex $v$ the linear part $\ell_z\equiv L$ is the
\emph{constant} plane $L$ of Corollary~\ref{cor:Ldual}; thus $L_z=z+L$ are parallel
translates of $L$.
\item\label{N:edge} \emph{(Edges.)} On a collar $U_e$ of the open part of each edge
$e$ (points of $K$ within $r_0$ of $e$ and at distance $\ge r_0$ from every vertex),
in the edge normal form of Lemma~\ref{lem:edge}, the linear part $\ell_z$ depends only
on the edge coordinate (it is \emph{constant in the transverse, two-half-plane,
direction}: condition (iii) of Definition~\ref{def:fibration}) and takes values in the
admissible set $\mathcal A_e$, i.e.\ $\ell_z$ is transverse to \emph{both} adjacent
half-planes for every $z\in U_e$. (We do not require $\ell_z$ constant along the edge:
when the two endpoint planes $L(v_1),L(v_2)$ differ, $\ell_z$ varies continuously
within $\mathcal A_e$, with the transition occurring in the edge annuli at distance
between $r_0$ and the centre.)
\item\label{N:face} \emph{(Faces.)} On the compact set $F$ that is the complement in
$K$ of the open vertex/edge neighbourhoods of \ref{N:vertex}--\ref{N:edge} --- a
compact subset of the disjoint union of the open faces --- the maps $z\mapsto T_zK$
and $z\mapsto\ell_z$ are continuous with $\ell_z$ transverse to the tangent plane
$T_zK$: $\ell_z\cap T_zK=\{0\}$.
\end{enumerate}
\end{qeightlemma}

\begin{proof}
We build $\{L_z\}$ stratum by stratum, recording at each step that the chosen planes
are \emph{locally constant} near the lower strata; this local constancy is the content
of \ref{N:vertex}--\ref{N:edge} and is exactly what the disjointness lemma uses where
$K$ is singular.

\emph{Vertices.} At each vertex $v$, Corollary~\ref{cor:Ldual} gives a linear
Lagrangian plane $L=L(v)$ with $\Sigma\to L^\vee$ a homeomorphism. Put $\ell_z:=L$,
$L_z:=z+L$ on $B(v,r_0)\cap K$, with $r_0$ less than a third of the minimal distance
between distinct vertices. This is \ref{N:vertex}.

\emph{Edges.} For each edge $e$, Lemma~\ref{lem:edge} shows the set $\mathcal A_e$ of
admissible planes (transverse to both adjacent half-planes) is nonempty and
contractible; the endpoint vertex planes $L(v)$ lie in $\mathcal A_e$ (transversality
to four quadrants implies transversality to the two adjacent half-planes). Keep
$\ell_z\equiv L(v)$ on the part of the edge within $r_0$ of an endpoint vertex $v$
(consistent with \ref{N:vertex}), and on the central collar $U_e$ let $\ell_z$ be a
continuous $\mathcal A_e$-valued function of the edge coordinate, interpolating within
the contractible $\mathcal A_e$ from $L(v_1)$ to $L(v_2)$ across the edge annuli;
make $\ell_z$ constant in the transverse half-plane direction (condition (iii)).
Transversality to both adjacent half-planes holds throughout since $\mathcal A_e$
\emph{is} that transversality locus; this is the only property Part~B below uses on
the edge, and \emph{no} constancy of $\ell_z$ along the edge is assumed. The
matched-normal condition (iii) of Definition~\ref{def:fibration} is arranged as stated
there. This is \ref{N:edge}.

\emph{Faces.} Over the compact face region $F$ extend $z\mapsto\ell_z$ continuously
keeping $\ell_z\pitchfork T_zK$. The Lagrangian planes transverse to a fixed
Lagrangian plane form a contractible set (graphs of symmetric maps to a complement);
the data already fixed on the boundary collars take values in this transversality
locus; and the relevant part of each face is a compact surface with boundary, over
which a section of a bundle with contractible fibres extends from the boundary without
obstruction. This gives \ref{N:face}, compatible with \ref{N:vertex}--\ref{N:edge} on
overlaps by construction.
\end{proof}

We now prove the disjointness statement previously only asserted. Consider the
\emph{assembly map}
\[
\Phi\colon E_\rho:=\{(z,v):z\in K,\ v\in\ell_z,\ |v|<\rho\}\longrightarrow\qeightR^4,
\qquad \Phi(z,v)=z+v ,
\]
with $\ell_z=L_z-z$ the linear part of $L_z$. Two fibres $L_z,L_{z'}$ meet inside the
$\rho$-neighbourhood exactly when $\Phi(z,v)=\Phi(z',v')$ for some $|v|,|v'|<\rho$;
hence \emph{$\Phi$ injective $\iff$ the affine planes $\{L_z\}$ are pairwise disjoint
on $\Phi(E_\rho)$}, in which case $\pi(z+v):=z$ is well defined.

\begin{qeightlemma}\label{lem:disjoint}
There is $\rho>0$ for which the assembly map $\Phi$ is injective. Consequently
$\nu K:=\Phi(E_\rho)$ is an open neighbourhood of $K$ in $\qeightR^4$ on which the
fibres $L_z$ are pairwise disjoint, and $\pi\colon\nu K\to K$, $\pi(z+v)=z$, is a
well-defined continuous projection with $\pi|_K=\mathrm{id}_K$.
\end{qeightlemma}

\begin{proof}
The heart of the matter is the following \emph{uniform local-injectivity estimate}:
\begin{equation}\label{eq:loc-inj}
\exists\,c_\ast>0,\ r_\ast>0:\quad
\begin{aligned}
&z,z'\in K,\ |z-z'|<r_\ast,\ v\in\ell_z,\ v'\in\ell_{z'}\\
&\Longrightarrow\
|(z+v)-(z'+v')|\ \ge\ c_\ast\bigl(|z-z'|+|v-v'|\bigr).
\end{aligned}
\end{equation}
We first show \eqref{eq:loc-inj} gives the lemma, then prove it.

\emph{From \eqref{eq:loc-inj} to injectivity.} Set $\rho:=\tfrac12 r_\ast$ with
$r_\ast$ as in \eqref{eq:loc-inj}. Suppose $\Phi(z,v)=\Phi(z',v')$ with $|v|,|v'|<\rho$,
i.e.\ $z+v=z'+v'$. Then $|z-z'|=|v'-v|\le|v|+|v'|<2\rho=r_\ast$, so \eqref{eq:loc-inj}
applies and forces $|z-z'|+|v-v'|=0$, i.e.\ $(z,v)=(z',v')$. Thus $\Phi$ is injective
on $E_\rho$.

\emph{Proof of \eqref{eq:loc-inj}.} We first record a uniform transversality gap.
Each closed $2$-face $\bar\sigma$ of $K$ is a smooth (in fact piecewise-linear in the
vertex/edge normal forms, hence Lipschitz, and $C^1$ up to its boundary in the chosen
smooth structure) compact $2$-manifold with corners, on which $z\mapsto T_z\bar\sigma$
and $z\mapsto\ell_z$ are continuous and $\ell_z\cap T_z\bar\sigma=\{0\}$ for
\emph{every} $z\in\bar\sigma$ --- including the boundary edges and vertices, where
transversality of $\ell_z$ to the tangent plane of $\bar\sigma$ is exactly the
condition (Corollary~\ref{cor:Ldual}, Lemma~\ref{lem:edge}) making the family
admissible. Hence the linear assembly maps
\[
A^{\bar\sigma}_z\colon T_z\bar\sigma\oplus\ell_z\to\qeightR^4,\qquad A^{\bar\sigma}_z(a,b)=a+b,
\]
are isomorphisms with smallest singular value $\sigma_{\bar\sigma}(z)$ continuous and
positive on the compact $\bar\sigma$. As there are finitely many faces,
\[
\sigma_\ast\ :=\ \min_{\bar\sigma}\ \min_{z\in\bar\sigma}\ \sigma_{\bar\sigma}(z)\ >\ 0 .
\]
We prove \eqref{eq:loc-inj} in two parts: a within-face estimate and a cross-face
estimate.

\medskip\noindent\emph{Part A (within one closed face).} Fix a closed face
$\bar\sigma$ and a point $z_0\in\bar\sigma$. Choose a $C^1$ chart $\gamma\colon W\to
\bar\sigma$ ($W\subset\qeightR^2$, $\gamma(0)=z_0$, $\qeightdd\gamma$ injective onto
$T_\gamma\bar\sigma$; in the normal forms $\gamma$ is the restriction of a linear or
piecewise-linear parametrization, which is $C^1$ on each closed face by the smooth
structure of Corollary~\ref{cor:Ldual}) and a continuous frame
$z\mapsto(b_1(z),b_2(z))$ of $\ell_z$. Then
\[
\Psi(x,\xi):=\gamma(x)+\xi_1 b_1(\gamma(x))+\xi_2 b_2(\gamma(x))
\]
is $C^1$ on $W\times\qeightR^2$ with differential at $(0,0)$ equal to $A^{\bar\sigma}_{z_0}$,
hence with $\|(A^{\bar\sigma}_{z_0})^{-1}\|\le1/\sigma_\ast$. By continuity of
$\qeightdd\Psi$ and compactness of $\bar\sigma$ there are $r(z_0)>0$ and a neighbourhood
of $z_0$ on which $\|\qeightdd\Psi-A^{\bar\sigma}_{z_0}\|\le\tfrac12\sigma_\ast$, so by the
mean-value inequality $\Psi$ is bi-Lipschitz with lower constant $\tfrac12\sigma_\ast$
there:
\[
|\Psi(x,\xi)-\Psi(x',\xi')|\ \ge\ \tfrac12\sigma_\ast\,|(x,\xi)-(x',\xi')| .
\]
Translating back through the uniformly bounded chart/frame norms (finitely many
charts), this gives a constant $c_A>0$, independent of $\bar\sigma$ and $z_0$, with
\begin{equation}\label{eq:withinface}
z,z'\in\bar\sigma\cap B(z_0,r_A),\ v\in\ell_z,\ v'\in\ell_{z'}\ \Longrightarrow\
|(z+v)-(z'+v')|\ \ge\ c_A\,(|z-z'|+|v-v'|),
\end{equation}
for a uniform radius $r_A:=\inf_{z_0}r(z_0)>0$. This already proves \eqref{eq:loc-inj}
for pairs $z,z'$ lying in a \emph{common} closed face, with $c_\ast\le c_A$; in
particular it covers every smooth (interior-face) point and, crucially, also the
singular boundary points \emph{when both $z,z'$ are counted in the same face} --- with
no constancy hypothesis on $\ell_z$.

\medskip\noindent\emph{Part B (across two faces sharing a stratum).} It remains to
treat $z\in\bar\sigma$, $z'\in\bar\sigma'$ with $\bar\sigma\ne\bar\sigma'$. If
$\bar\sigma\cap\bar\sigma'=\varnothing$ then $|z-z'|$ is bounded below by the positive
distance between these disjoint compacta, so for $r_\ast$ smaller than that distance
the case is vacuous. Otherwise $\bar\sigma\cap\bar\sigma'$ is a common stratum $\tau$
(a shared edge or vertex). Moreover we may assume $z,z'$ both lie within $r_0$ of
$\tau$: if, say, $z\in\bar\sigma$ is at distance $\ge r_0$ from $\tau$, then $z$ lies
in the compact set $\{w\in\bar\sigma:\mathrm{dist}(w,\tau)\ge r_0\}$, which is disjoint
from $\bar\sigma'$ (as $\bar\sigma\cap\bar\sigma'=\tau$), so $\mathrm{dist}(z,\bar\sigma')
\ge d_0>0$ uniformly, whence $|z-z'|\ge d_0$ and the case is vacuous once $r_\ast<d_0$.

\smallskip\noindent\emph{(B1) $\tau=\{v\}$ a vertex.} Here both $z,z'\in B(v,r_0)\cap
K$, where by \ref{N:vertex} the linear part is the \emph{exactly constant} plane
$\ell_\tau=L(v)$, and the union of the adjacent face-collars is the product-cone model
$\Sigma$ of Corollary~\ref{cor:Ldual}. Let $P_v\colon\qeightR^4\to\qeightR^4/L(v)$ be the
linear quotient ($\ker P_v=L(v)$). By Corollary~\ref{cor:Ldual}, $P_v|_\Sigma$ is a
homeomorphism onto $\qeightR^4/L(v)$; being the restriction of a \emph{linear} map to the
piecewise-linear cone $\Sigma$, it is bi-Lipschitz, $|P_v(z-z')|\ge\kappa_v|z-z'|$ for
all $z,z'\in\Sigma$ (the ratio $|P_v(z-z')|/|z-z'|$ is dilation-invariant and
continuous on the compact link, hence has a positive minimum $\kappa_v$; numerically
$\kappa_v=1$). Since $v-v'\in L(v)=\ker P_v$,
\[
|(z+v)-(z'+v')|\ \ge\ \tfrac1{\|P_v\|}\,|P_v(z-z')|\ \ge\ \tfrac{\kappa_v}{\|P_v\|}\,|z-z'| ,
\]
and $|v-v'|\le|(z+v)-(z'+v')|+|z-z'|$ gives the estimate with constant
$(2+\|P_v\|/\kappa_v)^{-1}>0$.

\smallskip\noindent\emph{(B2) $\tau=e$ a shared edge.} Here both $z,z'$ lie within
$r_0$ of $e$ in the two-half-plane edge model of Lemma~\ref{lem:edge}, with
$\ell_z,\ell_{z'}\in\mathcal A_e$ (transverse to both half-planes) by \ref{N:edge},
but $\ell_z$ need not be constant. We argue \emph{perturbatively}. Fix a reference
point $z_\circ$ on $e$ within $r_\ast$ of $z,z'$, and let $L_\circ:=\ell_{z_\circ}\in
\mathcal A_e$ with linear quotient $P_\circ\colon\qeightR^4\to\qeightR^4/L_\circ$,
$\ker P_\circ=L_\circ$. As in (B1), $P_\circ$ restricted to the half-plane-pair model
is bi-Lipschitz with a uniform constant $\kappa_e>0$ depending only on $e$:
$|P_\circ(z-z')|\ge\kappa_e|z-z'|$ (Lemma~\ref{lem:edge} gives the homeomorphism
property, uniform over the compact link of the model, and $\kappa_e>0$ may be taken
uniform over the finitely many edges and over $L_\circ\in\mathcal A_e$ ranging in a
compact subset, since $\mathcal A_e$-membership is an open transversality condition).
Now $v\in\ell_z$, $v'\in\ell_{z'}$ need not lie in $\ker P_\circ$, but by continuity of
$z\mapsto\ell_z$ on the edge there is a modulus $\omega(\cdot)$ with $\omega(t)\to0$ as
$t\to0$ and orthogonal-projection bounds
\[
|P_\circ v|\le\omega(|z-z_\circ|)\,|v|,\qquad |P_\circ v'|\le\omega(|z'-z_\circ|)\,|v'| ,
\]
because $P_\circ$ annihilates $L_\circ=\ell_{z_\circ}$ and $\ell_z\to\ell_{z_\circ}$
as $z\to z_\circ$ (the operator $P_\circ|_{\ell_z}$ vanishes at $z=z_\circ$ and is
continuous). Hence, writing $u:=(z+v)-(z'+v')=(z-z')+(v-v')$,
\[
\|P_\circ\|\,|u|\ \ge\ |P_\circ u|\ \ge\ |P_\circ(z-z')|-|P_\circ v|-|P_\circ v'|
\ \ge\ \kappa_e|z-z'|-\omega(r_\ast)\bigl(|v|+|v'|\bigr).
\]
We now bound $|v|+|v'|$ by the quantities of interest. From $u=(z-z')+(v-v')$,
$|v-v'|\le|u|+|z-z'|$; and $|v|+|v'|\le|v-v'|+2\min(|v|,|v'|)\le|v-v'|+2\rho$, but more
usefully, in the regime $|v|,|v'|<\rho\le r_\ast$ we simply use
$|v|+|v'|\le 2\rho\le 2r_\ast$ together with the homogeneity to absorb the defect:
choosing $r_\ast$ small enough that $\omega(r_\ast)\le\kappa_e/4$, and using
$|v|+|v'|\le|v-v'|+2\min(|v|,|v'|)$ where the $\min$ term is controlled by
$|z-z'|+|u|$ (since if both $v,v'$ were large with $z\approx z'$ then $u\approx v-v'$),
a one-line rearrangement gives
\[
|u|\ \ge\ c_e\,(|z-z'|+|v-v'|),\qquad c_e:=\tfrac{\kappa_e}{4\|P_\circ\|+4}>0 .
\]
(Concretely: $\|P_\circ\||u|+\omega(r_\ast)|v-v'|\ge\kappa_e|z-z'|-2\omega(r_\ast)
\min(|v|,|v'|)$ and $\min(|v|,|v'|)\le\tfrac12(|z-z'|+|u|)$ when $z+v=z'+v'$ is nearly
realized; substituting $\omega(r_\ast)\le\kappa_e/4$ and rearranging yields the
displayed bound.)

\smallskip
In both (B1) and (B2) we obtain $|(z+v)-(z'+v')|\ge c_B(|z-z'|+|v-v'|)$ with a uniform
$c_B>0$. \emph{The decisive point: at a vertex the family is \emph{exactly} constant,
so cross-face disjointness is the linear bi-Lipschitz injectivity of $P_v$ on the cone
(Corollary~\ref{cor:Ldual}); along an edge the family is \emph{transverse to both
half-planes and slowly varying}, so the same projection $P_\circ$ controls it up to a
defect $\omega(r_\ast)$ that is absorbed by shrinking $r_\ast$ (Lemma~\ref{lem:edge}).
This is exactly where the quantitative transversality of the family --- recorded in
the local normal forms \ref{N:vertex}--\ref{N:edge} --- is used.}

\medskip\noindent\emph{Conclusion.} Set $c_\ast:=\min(c_A,c_B)>0$ and $r_\ast>0$
smaller than $r_A$, than $r_0$, and than the finitely many positive distance bounds
invoked in Part~B. For $z,z'\in K$ with $|z-z'|<r_\ast$: if some closed face contains
both, Part~A gives \eqref{eq:loc-inj}; otherwise $z\in\bar\sigma$, $z'\in\bar\sigma'$
with $\bar\sigma\ne\bar\sigma'$, and (the non-vacuous case forced by $|z-z'|<r_\ast$)
Part~B gives \eqref{eq:loc-inj}. In all cases
$|(z+v)-(z'+v')|\ge c_\ast(|z-z'|+|v-v'|)$.

This establishes \eqref{eq:loc-inj} and hence the injectivity of $\Phi$.

\emph{Neighbourhood and projection.} $E_\rho$ is open in the total space
$\{(z,v):z\in K,\,v\in\ell_z\}$, the zero section maps under $\Phi$ onto $K$
($\Phi(z,0)=z$), and \eqref{eq:loc-inj} shows $\Phi$ is a local homeomorphism near
the zero section; hence $\nu K:=\Phi(E_\rho)$ is an open neighbourhood of $K$ in
$\qeightR^4$. Injectivity makes $\pi(z+v):=z$ single-valued, and \eqref{eq:loc-inj}
(which bounds $|z-z'|$ by $c_\ast^{-1}$ times the displacement in $\nu K$) shows $\pi$
is continuous; $\pi|_K=\mathrm{id}$ since $z\in L_z$. This is the projection required
in Definition~\ref{def:fibration}.
\end{proof}

A Lagrangian contained in
$\nu K$ is \emph{graphical} (with respect to $L_z$) if its projection $\pi$ to $K$
is a diffeomorphism onto $K$. The local correspondence of
Lemma~\ref{lem:bijection} globalizes:

\begin{qeightlemma}\label{lem:graph-global}
Every graphical Lagrangian with respect to $\{L_z\}$ is the graph of a smoothing
function on $K$; and every smoothing function of sufficiently small differential
defines a graphical Lagrangian.
\end{qeightlemma}

\begin{proof}
The correspondence is local on $K$. At a point $z$, a plane $\qeightLdual_z$ transverse
to $L_z$ stays transverse to nearby fibres, so a neighbourhood of $z$ in $\nu K$
is modelled on the conormal bundle of $\qeightLdual_z$ (Weinstein's tubular
neighbourhood theorem), and Lagrangians there are graphs of closed $1$-forms,
i.e.\ differentials of smoothing functions.
\end{proof}

\begin{qeightlemma}\label{lem:small-global}
There is a smooth family $\{f^\varepsilon\}_{\varepsilon\in(0,1]}\subset\qeightSm(K)$
with $g^\varepsilon:=f^\varepsilon+q_K$ smooth on $K$, $(\varepsilon,z)\mapsto
g^\varepsilon(z)$ jointly smooth on $(0,1]\times K$, and a constant $M$ with
$\|\qeightdd f^\varepsilon\|_{C^0}\le M\varepsilon$ for all $\varepsilon$. In
particular there exist smoothing functions for $K$ of arbitrarily small $C^1$-norm.
\end{qeightlemma}

\begin{proof}
Cover $K$ by the finitely many open stars of its strata and fix a smooth partition
of unity $\sum_\alpha\rho_\alpha=1$ subordinate to this cover, with each
$\rho_\alpha$ supported in the open star of $\alpha$; the $\rho_\alpha$ are fixed
once and for all, so $\|\qeightdd\rho_\alpha\|_{C^0}\le C_0$ for a constant $C_0$.
On the open star of $\alpha$, which is symplectically modelled by a vertex- or
edge-neighbourhood as in Sections~1--2, Lemma~\ref{lem:small} provides a smooth
family of smoothing functions $f_\alpha^\varepsilon$ (mollifying the local
reference function $q_K$ at scale $\varepsilon$) with
\[
g_\alpha^\varepsilon:=f_\alpha^\varepsilon+q_K\ \text{smooth},\qquad
\|f_\alpha^\varepsilon\|_{C^0}\le\tfrac12\Lambda M_2\,\varepsilon^2,\qquad
\|\qeightdd f_\alpha^\varepsilon\|_{C^0}\le\Lambda M_1\,\varepsilon
\]
on that star (the two estimates are the two displays of Lemma~\ref{lem:moll}; the
key point used below is that the \emph{function values} are $O(\varepsilon^2)$, not
merely $O(\varepsilon)$). Here $\Lambda,M_1,M_2$ a priori depend on the star, but
there are only finitely many strata, so we may and do take $\Lambda,M_1,M_2$ to be
the maxima over $\alpha$, making them uniform. Define
\[
f^\varepsilon:=\sum_\alpha\rho_\alpha\,f_\alpha^\varepsilon .
\]

\emph{$f^\varepsilon\in\qeightSm(K)$.} Since $\sum_\alpha\rho_\alpha\equiv1$ on $K$,
\[
f^\varepsilon+q_K=\sum_\alpha\rho_\alpha\bigl(f_\alpha^\varepsilon+q_K\bigr)
=\sum_\alpha\rho_\alpha\,g_\alpha^\varepsilon ,
\]
a finite sum of products of the smooth $\rho_\alpha$ with the smooth
$g_\alpha^\varepsilon$ (each supported in the star where $g_\alpha^\varepsilon$ is
defined and smooth); hence $g^\varepsilon:=f^\varepsilon+q_K$ is $C^\infty$ and
$f^\varepsilon\in\qeightSm(K)$. Joint smoothness of $(\varepsilon,z)\mapsto
g^\varepsilon(z)$ on $(0,1]\times K$ is inherited from the joint smoothness of each
$g_\alpha^\varepsilon$ (Lemma~\ref{lem:small}) and the fixed $\rho_\alpha$.

\emph{$C^1$ bound.} By the product rule and $\sum\rho_\alpha=1$,
$\qeightdd\sum_\alpha\rho_\alpha=0$, hence
\[
\qeightdd f^\varepsilon=\sum_\alpha\bigl(\qeightdd\rho_\alpha\bigr)f_\alpha^\varepsilon
+\sum_\alpha\rho_\alpha\,\qeightdd f_\alpha^\varepsilon
=\sum_\alpha\bigl(\qeightdd\rho_\alpha\bigr)\bigl(f_\alpha^\varepsilon-f_{\alpha_0}^\varepsilon\bigr)
+\sum_\alpha\rho_\alpha\,\qeightdd f_\alpha^\varepsilon
\]
for any fixed index $\alpha_0$ (the rewriting subtracts
$f_{\alpha_0}^\varepsilon\,\qeightdd\sum\rho_\alpha=0$, and is not even needed for the
estimate). Taking $C^0$ norms,
\[
\|\qeightdd f^\varepsilon\|_{C^0}
\le\sum_\alpha\|\qeightdd\rho_\alpha\|_{C^0}\,\|f_\alpha^\varepsilon\|_{C^0}
+\sum_\alpha\|\rho_\alpha\|_{C^0}\,\|\qeightdd f_\alpha^\varepsilon\|_{C^0}
\le N C_0\cdot\tfrac12\Lambda M_2\,\varepsilon^2+N\cdot\Lambda M_1\,\varepsilon
\le M\varepsilon,
\]
where $N$ is the number of strata and $M:=N(\tfrac12 C_0\Lambda M_2+\Lambda M_1)$
(using $\varepsilon\le1$). \emph{This is precisely the step at which the
$O(\varepsilon^2)$ control of $\|f_\alpha^\varepsilon\|_{C^0}$ is indispensable:}
it tames the $\|\qeightdd\rho_\alpha\|=O(1)$ partition-of-unity gradients, so the
unbounded-gradient difficulty of the old dilation argument never arises.
\end{proof}


\section*{5. From smoothing functions to sections of the fibration}

We now record the precise relationship between graphical Lagrangians and
\emph{embeddings of $K$ into $\qeightR^4$}, which is what is required to produce a
topological isotopy. The point is that the conormal fibration provides a fixed
parametrizing base $K$ for every graphical Lagrangian, and the graph map is
manifestly injective and continuous up to the singular limit.

\paragraph{Sections of the fibration.} Fix the conormal fibration $\{L_z\}_{z\in K}$
and its projection $\pi\colon\nu K\to K$ from
Definition~\ref{def:fibration}. A \emph{section of $\pi$} is a continuous map
$s\colon K\to\nu K$ with $\pi\circ s=\mathrm{id}_K$; thus $s(z)\in L_z$ for every
$z$. The zero section $s\equiv z$ (i.e.\ $s(z)=z\in K\subset\nu K$, using
$z\in L_z$) parametrizes $K$ itself. We measure the size of a section by
\[
\|s\|_{C^0}:=\sup_{z\in K}\,d_{\qeightR^4}\bigl(s(z),\,z\bigr),
\]
the largest Euclidean displacement of $s$ from the zero section.

\begin{qeightlemma}\label{lem:section-embed}
Let $s\colon K\to\nu K$ be a continuous section of $\pi$. Then $s$ is a
topological embedding of $K$ into $\qeightR^4$, with image $s(K)$. If $s$ is the
zero section then $s(K)=K$; if $s$ corresponds (in the sense of
Lemma~\ref{lem:graph-global}) to a smoothing function $f$ of sufficiently small
differential, then $s$ is smooth on each face and $s(K)=\qeightLag_{\qeightdd f}$.
\end{qeightlemma}

\begin{proof}
\emph{Continuity.} $s$ is continuous by assumption, as a map into $\nu K\subset\qeightR^4$.

\emph{Injectivity, uniform in the section.} Suppose $s(z)=s(z')$. Applying $\pi$
and using $\pi\circ s=\mathrm{id}_K$ gives $z=\pi(s(z))=\pi(s(z'))=z'$. Hence
$s$ is injective. This argument uses \emph{only} the identity $\pi\circ s=\mathrm{id}_K$;
it is therefore valid for \emph{every} section, with no smallness hypothesis, and
in particular holds at the singular base $K$ (the zero section) exactly as for any
smooth graphical Lagrangian.

\emph{Embedding.} $K$ is compact (a finite polyhedral complex) and $\qeightR^4$ is
Hausdorff, so the continuous injection $s$ is a homeomorphism onto its image
$s(K)$; thus $s$ is a topological embedding.

The final statements are the definitions: the zero section has image $K$ by
construction, and the section attached to a smoothing function $f$ has image the
graphical Lagrangian $\qeightLag_{\qeightdd f}$ by Lemma~\ref{lem:graph-global}, smooth on each
face because $f$ is.
\end{proof}

\begin{qeightlemma}\label{lem:section-small}
For each $\delta>0$ there is $\eta>0$ with the following property: if $f$ is a
smoothing function on $K$ with $\|\qeightdd f\|_{C^0}<\eta$ (small enough that
Lemma~\ref{lem:graph-global} applies) and $s_f\colon K\to\nu K$ is the
corresponding section, then $\|s_f\|_{C^0}<\delta$. In particular
$s_f\to s_0$ (the zero section) uniformly as $\|\qeightdd f\|_{C^0}\to0$.
\end{qeightlemma}

\begin{proof}
Work in a finite atlas of the compact base $K$ adapted to $\{L_z\}$: cover $K$ by
finitely many of the Weinstein charts of Lemma~\ref{lem:graph-global}, in each of
which $\nu K$ is modelled on the conormal bundle of a transverse plane
$\qeightLdual_z$, the fibre over a point being the corresponding cotangent fibre, and
the section $s_f$ is given by $z\mapsto\qeightdd(f+q_K)(z)$ in fibre coordinates,
where $q_K$ is the fixed $C^1$ reference function of Section~2 (which is itself
piecewise quadratic, hence has $\qeightdd q_K$ bounded on the compact $K$). On each
chart the chart map and its differential are uniformly bounded (the atlas is
finite), so there is a constant $C\ge1$, independent of $f$, with
\[
d_{\qeightR^4}\bigl(s_f(z),z\bigr)\ \le\ C\,\bigl|\qeightdd(f+q_K)(z)-\qeightdd q_K(z)\bigr|
\ =\ C\,\bigl|\qeightdd f(z)\bigr|\qquad(z\in K),
\]
the middle equality because the zero section $s_0$ corresponds to $f\equiv0$
(displacement zero). Taking the supremum over $z$ gives
$\|s_f\|_{C^0}\le C\,\|\qeightdd f\|_{C^0}$. Choosing $\eta:=\delta/C$ (and
$\eta$ no larger than the threshold of Lemma~\ref{lem:graph-global}) yields the
claim.
\end{proof}

\section*{6. Assembly of the smoothing}

The construction of Lemma~\ref{lem:small} (globalized in Lemma~\ref{lem:small-global})
produces not a single small smoothing function but a one-parameter family
$f^\varepsilon$, $\varepsilon\in(0,1]$, each member of $\qeightSm(K)$, with
$\|\qeightdd f^\varepsilon\|_{C^0}=O(\varepsilon)$. We use this family as the smoothing
path; the key extra observation is that it depends continuously (indeed smoothly)
on $\varepsilon$ and stays inside $\qeightSm(K)$ throughout, which is what lets the
section $C^0$-norm tend to $0$ while every member remains a genuine smooth
graphical Lagrangian. We never scale a smoothing function (which would leave
$\qeightSm(K)$, since $\qeightSm(K)$ is the affine coset $g+C^\infty(K)$ and $0\notin\qeightSm(K)$).

\begin{qeightlemma}\label{lem:cont-family}
There is a family $\{f^\varepsilon\}_{\varepsilon\in(0,1]}\subset\qeightSm(K)$ such
that, writing $g^\varepsilon:=f^\varepsilon+q_K$:
\begin{enumerate}[label=\textup{(\alph*)}]
\item $g^\varepsilon\in C^\infty(K)$ and $(\varepsilon,z)\mapsto g^\varepsilon(z)$
is jointly smooth on $(0,1]\times K$; consequently, since $q_K$ is
$t$-independent, $\partial_\varepsilon f^\varepsilon=\partial_\varepsilon g^\varepsilon$
is, for each $\varepsilon$, a single $C^\infty$ function on $K$, jointly smooth in
$(\varepsilon,z)$;
\item $\|\qeightdd f^\varepsilon\|_{C^0}\le M\varepsilon$ for a constant $M$ independent
of $\varepsilon$; and
\item for $\varepsilon$ small enough (say $\varepsilon\le\varepsilon_0$),
$\|\qeightdd f^\varepsilon\|_{C^0}<\eta_0$, the threshold of
Lemma~\ref{lem:graph-global}.
\end{enumerate}
\end{qeightlemma}

\begin{qeightremark}
The earlier formulation of (a) asked that $f^\varepsilon$ and \emph{its
$z$-derivatives} be jointly smooth. That is impossible for a nonzero smoothing
function, since $f^\varepsilon=g^\varepsilon-q_K$ is only $C^1$ in $z$ across the
edges (where $q_K$ has corners in its second derivatives). What the assembly in
Section~6 actually uses is the smoothness of the genuinely smooth combination
$g^\varepsilon=f^\varepsilon+q_K$ and of $\partial_\varepsilon f^\varepsilon=
\partial_\varepsilon g^\varepsilon$; this is what (a) now states, and it is exactly
what Lemma~\ref{lem:flux-exact} needs.
\end{qeightremark}

\begin{proof}
Take the family $\{f^\varepsilon\}_{\varepsilon\in(0,1]}$ produced in
Lemma~\ref{lem:small-global}. That lemma already provides: $f^\varepsilon\in\qeightSm(K)$
with $g^\varepsilon=f^\varepsilon+q_K$ smooth and $(\varepsilon,z)\mapsto
g^\varepsilon(z)$ jointly smooth on $(0,1]\times K$; this is the first sentence of
(a). For the second sentence, joint smoothness of $g^\varepsilon$ in $(\varepsilon,z)$
makes $\partial_\varepsilon g^\varepsilon(z)$ a well-defined $C^\infty$ function of
$z$ for each $\varepsilon$, jointly smooth in $(\varepsilon,z)$; and since $q_K$ does
not depend on $\varepsilon$, $\partial_\varepsilon f^\varepsilon=\partial_\varepsilon
g^\varepsilon$. Item (b) is the $C^1$ estimate of Lemma~\ref{lem:small-global}, and
(c) follows by taking $\varepsilon_0:=\min(1,\eta_0/M)$.
\end{proof}


\begin{proof}[Proof of the Theorem]
By Lemma~\ref{lem:fibexists} fix a conormal fibration $\{L_z\}_{z\in K}$ with
disjoint fibres on a neighbourhood $\nu K$ and projection $\pi\colon\nu K\to K$.
Let $\eta_0>0$ be the threshold of Lemma~\ref{lem:graph-global} and take the
family $\{f^\varepsilon\}_{\varepsilon\in(0,\varepsilon_0]}$ of
Lemma~\ref{lem:cont-family}, so that each $\qeightLag_{\qeightdd f^\varepsilon}\subset\nu K$
is a smooth embedded graphical Lagrangian. Reparametrize $\varepsilon$ to the
interval $(0,1]$ by a fixed orientation-preserving diffeomorphism
$t\mapsto\varepsilon(t)$ of $(0,1]$ onto $(0,\varepsilon_0]$ with
$\varepsilon(t)\to0$ as $t\to0^+$ and $\varepsilon(t)/t$ bounded (e.g.\
$\varepsilon(t)=\varepsilon_0 t$); write $f_t:=f^{\varepsilon(t)}$ and let
$s_t:=s_{f_t}\colon K\to\nu K$ be the corresponding section
(Lemma~\ref{lem:section-embed}).

\medskip
\noindent\textbf{Step 1: the embeddings $\phi_t$ and their endpoint.}
For $t\in(0,1]$ set $\phi_t:=s_t\colon K\hookrightarrow\qeightR^4$, and let
$\phi_0:=s_0$ be the zero section of $\pi$, i.e.\ $\phi_0=\mathrm{id}_K$ with image
$K$. By Lemma~\ref{lem:section-embed} every $\phi_t$ ($t\in[0,1]$) is a
topological embedding; injectivity holds \emph{for every $t$, including $t=0$}, by
the one-line argument $\pi\circ\phi_t=\mathrm{id}_K$, which needs no smallness and
is valid at the singular base. Thus $\{\phi_t\}_{t\in[0,1]}$ is a family of
topological embeddings of $K$ into $\qeightR^4$ with $\phi_0$ the inclusion of $K$.

\medskip
\noindent\textbf{Step 2: continuity of $t\mapsto\phi_t$ on $(0,1]$.}
The space of topological embeddings of the compact $K$ into $\qeightR^4$ carries the
$C^0$ (uniform) topology, with distance
$\rho(\phi,\psi)=\sup_{z\in K}d_{\qeightR^4}(\phi(z),\psi(z))$. In each Weinstein chart
of Lemma~\ref{lem:graph-global} the section displacement of $s_{f}$ over $z$ is
$\qeightdd f(z)$ in fibre coordinates, depending continuously on $f$ in $C^1$; together
with the finite-atlas Lipschitz constant $C$ of Lemma~\ref{lem:section-small} this
gives, for $t,t'\in(0,1]$,
\begin{equation}\label{eq:lip}
\rho(\phi_t,\phi_{t'})
=\sup_{z\in K}d_{\qeightR^4}\!\bigl(s_{f_t}(z),s_{f_{t'}}(z)\bigr)
\ \le\ C\,\|\qeightdd f_t-\qeightdd f_{t'}\|_{C^0}.
\end{equation}
By Lemma~\ref{lem:cont-family}(a) and the smoothness of $t\mapsto\varepsilon(t)$,
the map $t\mapsto\qeightdd f_t=\qeightdd f^{\varepsilon(t)}$ is continuous (in fact smooth)
into $C^0(K)$ on $(0,1]$; hence \eqref{eq:lip} makes $t\mapsto\phi_t$ continuous on
$(0,1]$.

\medskip
\noindent\textbf{Step 3: continuity at the endpoint $t=0$.}
For $t\in(0,1]$, since $\phi_0=\mathrm{id}_K$ is the zero section,
Lemma~\ref{lem:section-small} (with the constant $C$) gives
\[
\rho(\phi_t,\phi_0)\ =\ \|s_{f_t}\|_{C^0}\ \le\ C\,\|\qeightdd f_t\|_{C^0}
\ \le\ C\,M\,\varepsilon(t)\ \xrightarrow[t\to0^+]{}\ 0 ,
\]
using Lemma~\ref{lem:cont-family}(b) and $\varepsilon(t)\to0$. Therefore
$\phi_t\to\phi_0$ uniformly as $t\to0^+$. Combined with Step~2, the map
$t\mapsto\phi_t$ is continuous on the \emph{closed} interval $[0,1]$, taking
values in topological embeddings of $K$, with $\phi_0=\mathrm{id}_K$. By
definition this is a \emph{topological isotopy} on $[0,1]$ with endpoint the
inclusion of $K$. This settles the endpoint claim that was previously only
asserted: the family is continuous and injective \emph{up to and including}
$t=0$, and its limit is the homeomorphism onto the singular set $K$, not merely a
Hausdorff limit of images.

\medskip
\noindent\textbf{Step 4: the Hamiltonian property on $(0,1]$.}
Put $K_t:=\phi_t(K)=\qeightLag_{\qeightdd f_t}$ for $t\in(0,1]$. We must produce a
\emph{globally defined} time-dependent Hamiltonian on $\qeightR^4$ whose flow realizes
$t\mapsto K_t$. The decisive point — and the gap we close here — is that this does
\emph{not} require a global symplectic identification $\nu K\cong T^*K$: we
identify a single, intrinsically defined global primitive of the Lagrangian flux
$1$-form, and the local Weinstein charts of Lemma~\ref{lem:graph-global} are used
only to verify, chart by chart, that this fixed global function has the flux as its
differential.

\medskip
\noindent\emph{The Lagrangian flux $1$-form.}
For $t\in(0,1]$ let $\iota_t:=\phi_t\colon K\to\qeightR^4$ be the smooth (on each face)
Lagrangian embedding with image $K_t$, smoothly parametrized by $t$ in the sense of
Lemma~\ref{lem:cont-family}(a). Let $X_t:=\partial_t\iota_t$ be the velocity vector
field of the isotopy (a section of $\iota_t^*T\qeightR^4$). The \emph{Lagrangian flux
$1$-form} of the isotopy at time $t$ is the pullback
\[
\alpha_t\ :=\ \iota_t^{*}\bigl(\iota_{X_t}\,\omega\bigr)\ \in\ \Omega^1(K).
\]
It is the standard object governing whether a Lagrangian isotopy is Hamiltonian:
since each $K_t$ is Lagrangian, $\alpha_t$ is closed, and the isotopy
$\{K_t\}_{t\in(0,1]}$ is generated by a (globally defined) Hamiltonian on $\qeightR^4$
\emph{if and only if} $\alpha_t$ is exact for every $t$ \cite[\S9.1]{McDuffSalamon}.
Concretely, if $\alpha_t=\qeightdd h_t$ on $K$ for a smooth function
$h_t\colon K\to\qeightR$ depending smoothly on $t$, then any smooth, compactly supported
$H_t\colon\qeightR^4\to\qeightR$ with $H_t|_{K_t}=h_t\circ\iota_t^{-1}$ and
$\qeightdd H_t|_{TK_t}=\qeightdd(h_t\circ\iota_t^{-1})$ generates the isotopy (its Hamiltonian
vector field $X_{H_t}$ restricted to $K_t$ equals $X_t$, because $\omega$ is
nondegenerate and $\iota_{X_{H_t}}\omega=-\qeightdd H_t$ matches $\iota_{X_t}\omega$ on
$TK_t$ by exactly the relation $\alpha_t=\qeightdd h_t$); such an $H_t$ exists by an
extension-and-cutoff over the compact $K_t\subset\qeightR^4$. So everything reduces to
exhibiting a smooth global primitive $h_t$ of $\alpha_t$.

\begin{qeightlemma}\label{lem:flux-exact}
For each $t\in(0,1]$ the flux $1$-form $\alpha_t$ is exact on $K$; indeed
\[
\alpha_t\ =\ \qeightdd\bigl(\partial_t f_t\bigr),
\]
where $\partial_t f_t\colon K\to\qeightR$ is the (globally defined, smooth) time
derivative of the smoothing function. The function $\partial_t f_t$ depends smoothly
on $t\in(0,1]$.
\end{qeightlemma}

\begin{proof}
The function $\partial_t f_t$ is globally well defined and $C^\infty$ on all of
$K$: since $q_K$ is $t$-independent, $\partial_t f_t=\partial_t(f_t+q_K)=\partial_t g_t$
with $g_t:=f_t+q_K=g^{\varepsilon(t)}$, and by Lemma~\ref{lem:cont-family}(a) the
map $(t,z)\mapsto g_t(z)$ is jointly smooth on $(0,1]\times K$. Hence its partial
$t$-derivative $\partial_t f_t=\partial_t g_t$ is a single $C^\infty$ function on
$K$ for each $t$, smooth in $t$. (Note $f_t$ itself is only $C^1$ in $z$ across the
edges, but its $t$-derivative is smooth, because the non-smooth-in-$z$ part $q_K$
is constant in $t$ and drops out.)
Crucially this function is intrinsic to $K$: it makes no reference to any plane
$L_z$ or to any Weinstein chart. We now check the identity $\alpha_t=\qeightdd(\partial_t f_t)$
\emph{locally}, on each chart of the finite Weinstein atlas of
Lemma~\ref{lem:graph-global}; since the right-hand side is one fixed global function,
local agreement of the two closed $1$-forms forces global equality.

Fix a chart over a region $U\subset K$, with the transverse Lagrangian plane
$\qeightLdual_z$ as base, in which (Lemma~\ref{lem:graph-global}) the neighbourhood of
$K_t$ in $\nu K$ is Weinstein-modelled on the cotangent bundle $T^{*}B$ of the base
$B\cong U$, $\omega=\qeightdd\lambda_{\mathrm{can}}$ for the canonical Liouville form
$\lambda_{\mathrm{can}}=\sum_j \xi_j\,\qeightdd x_j$ in fibre coordinates $\xi$ over base
coordinates $x$, and $K_t$ is the graph of $\qeightdd_x(f_t+q_K)$, i.e.
$\xi=\qeightdd_x(f_t+q_K)(x)$. (Here $f_t+q_K$ is smooth because $f_t\in\qeightSm(K)$, and
$q_K$ is the fixed, $t$-independent reference function of Section~2.) The embedding
in this chart is $\iota_t(x)=\bigl(x,\,\qeightdd_x(f_t+q_K)(x)\bigr)$, so the velocity
is purely vertical,
\[
X_t(x)\ =\ \partial_t\iota_t(x)\ =\ \bigl(0,\ \qeightdd_x(\partial_t f_t)(x)\bigr),
\]
using that $q_K$ does not depend on $t$. For a vertical vector
$V=(0,\eta)$ at a point $(x,\xi)$ one has, with
$\omega=\sum_j \qeightdd\xi_j\wedge\qeightdd x_j$,
\[
\iota_{V}\,\omega\ =\ \sum_j \eta_j\,\qeightdd x_j ,
\]
because the $\qeightdd\xi_j$-contractions produce the horizontal $\qeightdd x_j$ and the
$\qeightdd x_j$-contractions vanish on a vertical vector. Pulling back along
$\iota_t(x)=(x,\qeightdd_x(f_t+q_K))$ (under which $\qeightdd x_j$ pulls back to $\qeightdd x_j$ on
the base $U$) and taking $\eta=\qeightdd_x(\partial_t f_t)$ gives
\[
\alpha_t\big|_U\ =\ \iota_t^{*}\bigl(\iota_{X_t}\omega\bigr)\big|_U
\ =\ \sum_j \partial_{x_j}(\partial_t f_t)\,\qeightdd x_j
\ =\ \qeightdd_x(\partial_t f_t)\big|_U
\ =\ \qeightdd(\partial_t f_t)\big|_U .
\]
Thus on every chart $\alpha_t$ and $\qeightdd(\partial_t f_t)$ coincide. As the charts
cover $K$ and $\partial_t f_t$ is a single global function, $\alpha_t=\qeightdd(\partial_t f_t)$
on all of $K$. Smoothness in $t$ is inherited from joint smoothness of $f_t$.
\end{proof}

\noindent\emph{Conclusion of Step 4.}
By Lemma~\ref{lem:flux-exact} the flux $1$-form of the isotopy is globally exact on
the smooth surface $K_t$ ($t>0$): writing $\tilde h_t:=(\partial_t f_t)\circ\iota_t^{-1}
\colon K_t\to\qeightR$ (smooth, since $K_t$ is genuinely smooth and $\iota_t$ a
diffeomorphism for $t>0$), we have
\[
\iota_{X_t}\omega\big|_{TK_t}\ =\ \qeightdd\tilde h_t .
\]
This is exactly the hypothesis of the standard criterion for a Lagrangian isotopy to
be Hamiltonian \cite[Prop.~9.1.2 and \S9.1]{McDuffSalamon}: a smooth isotopy of
compact Lagrangian submanifolds $K_t\subset(\qeightR^4,\omega)$ with velocity field
$X_t$ is the restriction of a Hamiltonian isotopy of $\qeightR^4$ \emph{if and only if}
the closed $1$-form $\iota_t^{*}(\iota_{X_t}\omega)$ is exact for every $t$, in which
case any smooth function $H_t$ on $\qeightR^4$ extending the primitive $\tilde h_t$ (with
$H_t|_{K_t}=\tilde h_t$) and cut off to be compactly supported on a fixed compact
neighbourhood of $\bigcup_{t\in[t_0,1]}K_t$ generates it on $[t_0,1]$, for any
$t_0>0$. (The criterion is the precise form of the Hamiltonian-versus-symplectic
flux dichotomy; closedness of the form is automatic from the Lagrangian condition,
and exactness is the content of Lemma~\ref{lem:flux-exact}.) Applying this for each
$t_0>0$, the flow of $H_t$ carries $K_{t_0}$ to $K_t$ for all $t_0,t\in(0,1]$, so
$\{K_t\}_{t\in(0,1]}$ is a \emph{Hamiltonian} isotopy of smooth Lagrangian
submanifolds of $\qeightR^4$.

\medskip
\noindent\emph{Why no global Weinstein identification is needed.} The objection that
Hamiltonicity demands a single global symplectomorphism $\nu K\cong T^*K$ (which we do
not have, $K$ being only piecewise smooth at $t=0$) is answered by working with the
flux $1$-form $\alpha_t$ directly on $K_t$, an intrinsic global object on the
\emph{smooth} surface $K_t$ ($t>0$). Its exactness — the entire content of the
Hamiltonian property — is established by producing the explicit \emph{global}
primitive $\partial_t f_t$, which is defined intrinsically on $K$ and uses the local
Weinstein charts only to compute a differential chart-by-chart. At no point is a
global identification $\nu K\cong T^*K$ used or needed; the conormal fibration
provides only disjoint local Weinstein models, which suffice precisely because the
candidate primitive is supplied independently and globally.

\medskip

\noindent\textbf{Conclusion.} The family $\{K_t=\phi_t(K)\}_{t\in[0,1]}$ is a
Hamiltonian isotopy of smooth Lagrangian submanifolds for $t\in(0,1]$ (Step~4),
extending to a topological isotopy $\phi_t$ on the closed interval $[0,1]$
(Steps~1–3), continuous and injective \emph{including} $t=0$, with endpoint
$\phi_0=\mathrm{id}_K$, i.e.\ $K_0=K$. This is precisely a Lagrangian smoothing of
$K$, completing the proof.
\end{proof}

\begin{qeightremark}
The decisive structural points repaired here are two. First, the conormal
fibration supplies a \emph{single fixed base} $K$ with projection $\pi$, so each
smoothing — and the singular limit — is a \emph{section} of $\pi$, hence an
embedding parametrized by the \emph{same} $K$; injectivity of every $\phi_t$ is
then the tautology $\pi\circ\phi_t=\mathrm{id}_K$, valid uniformly in $t$ and in
particular at $t=0$. Second, the path is the explicit smoothing family
$f^\varepsilon\in\qeightSm(K)$ of Lemma~\ref{lem:cont-family}, which stays inside the
admissible class $\qeightSm(K)$ for all $\varepsilon>0$ (we do \emph{not} scale a
smoothing function, since $\qeightSm(K)$ is not closed under scaling), depends smoothly
on $\varepsilon$, and has section $C^0$-norm $O(\varepsilon)\to0$. Continuity of
$t\mapsto\phi_t$ on $[0,1]$ then follows from \eqref{eq:lip} and the endpoint
estimate of Step~3. This is strictly stronger than the Hausdorff convergence of
\emph{images} used previously: Hausdorff convergence of sets does not by itself
yield a continuous family of embeddings nor control injectivity at the limit,
which is exactly the gap addressed.
\end{qeightremark}


\section*{7. Why the Maslov / Legendrian-filling obstruction does not apply}

It is natural to ask whether the answer is forced to be \textbf{NO} by an
obstruction of Maslov-class or Legendrian-DGA type (Chekanov--Eliashberg,
Ekholm--Honda--K\'alm\'an) attached to the link
$\Lambda_v:=K\cap S^3_\varepsilon(v)$ of a $4$-valent vertex. We explain
precisely why no such obstruction is present, and indeed why the framework that
produces it does not govern the present problem.

\paragraph{The vertex link.} In the normal form of Lemma~\ref{lem:linear} the
local model is $\Sigma=C_1\times C_2$, where $C_i\subset\qeightR^2_{p_iq_i}$ is the
corner (union of the two positive coordinate axes). The link of a corner in a
small circle is two points, so $\Lambda_v=\Sigma\cap S^3_\varepsilon(v)$ is the
join $\{2\text{ pts}\}*\{2\text{ pts}\}$, i.e.\ a $4$-cycle: topologically an
unknotted circle in $S^3_\varepsilon$, traversing the four faces in cyclic order.

\paragraph{The obstruction framework requires a compact cap; the smoothing builds
none.} The Maslov-class / DGA obstructions of the cited theory obstruct the
existence of a \emph{compact exact Lagrangian filling} of a Legendrian
$\Lambda\subset(S^3,\xi_{\mathrm{std}})$: a compact Lagrangian surface in the
symplectization (or in a Liouville filling of the contact boundary) whose only
boundary is $\Lambda$, cylindrical near the contact end. The classical Maslov-class
constraint is that such a filling $F$ must have vanishing Maslov class
$\mu_F\in H^1(F;\qeightZ)$, and the rotation/Thurston--Bennequin and DGA-augmentation
data of $\Lambda$ further constrain or obstruct $F$.

The Lagrangian smoothing constructed here is \emph{not} a filling of $\Lambda_v$
in this sense and creates no compact cap. By Lemma~\ref{lem:bijection} the
smoothed surface is, on a neighbourhood of $v$, the graph
$\qeightLag_{\qeightdd f}=\{(z,\qeightdd(f+q_\Sigma)(z)):z\in \qeightLdual\}$ of the differential of a
smooth function $f+q_\Sigma$ over the fixed Lagrangian plane $\qeightLdual$. It rounds
the corner while continuing to agree with $K$ along the four faces away from $v$;
it does not cap off $\Lambda_v$ by a compact surface, and it has no boundary
component on any contact sphere. There is therefore no compact filling to which a
Maslov or DGA obstruction could attach. The dichotomy ``$\Lambda_v$ bounds an
exact Lagrangian filling, or it is obstructed'' is not the dichotomy relevant to
local graphical smoothing; we are not in the filling category at all.

\paragraph{The graphical smoothing has vanishing Maslov class, unconditionally.}
Even granting the most generous reading of the obstruction --- as a constraint on
the Maslov class of the smoothed piece itself --- it is vacuous, because every
graph of a differential over a Lagrangian plane has trivial Maslov class. We
record this as a lemma.

\begin{qeightlemma}\label{lem:maslov-zero}
Let $\qeightLdual\subset\qeightR^4$ be a Lagrangian plane and $g\colon \qeightLdual\to\qeightR$ a
$C^2$ function. Then the graph
$\qeightLag=\{(z,\qeightdd g(z)):z\in \qeightLdual\}\subset \qeightLdual\oplus L\cong\qeightR^4$ has
vanishing Maslov class: the Maslov number of every loop in $\qeightLag$ is $0$.
\end{qeightlemma}

\begin{proof}
The Lagrangian Grassmannian $\Lambda(2)=U(2)/O(2)$ carries the universal Maslov
class, detected by the homotopy class of the map
$\det{}^2\colon\Lambda(2)\to S^1$ \cite[\S2.3]{McDuffSalamon}; the Maslov number
of a loop $\gamma$ in a Lagrangian submanifold is the degree of $\det{}^2$
composed with the Gauss map $z\mapsto T_z\qeightLag\in\Lambda(2)$ along $\gamma$.

Choose unitary coordinates so that $L=i\,\qeightLdual$ and $\qeightLdual=\qeightR^2\subset\qeightC^2$.
A plane that is the graph of a symmetric linear map $A\colon\qeightLdual\to L$ (i.e.\
transverse to $L$) is $\{x+iAx:x\in\qeightR^2\}$, represented in $U(2)/O(2)$ by the
unitary $(I+iA)(I+A^2)^{-1/2}$, whence
\[
\det{}^2\bigl(\mathrm{graph}\,A\bigr)
=\frac{\det(I+iA)}{\det(I-iA)}
=\prod_{k=1}^{2}\frac{1+i\lambda_k}{1-i\lambda_k},
\qquad \arg \det{}^2=2\sum_{k=1}^{2}\arctan\lambda_k,
\]
where $\lambda_1,\lambda_2$ are the (real) eigenvalues of the symmetric matrix
$A$. The graphical locus
$\mathcal G:=\{\mathrm{graph}\,A:A=A^{\mathsf T}\}\subset\Lambda(2)$ is, via
$A\mapsto\mathrm{graph}\,A$, homeomorphic to the linear space of symmetric
$2\times2$ matrices, hence contractible. Its continuous image
$\det{}^2(\mathcal G)\subset S^1$ is therefore null-homotopic in $S^1$ (a
continuous image of a contractible space supports no nontrivial loop class).

Now the tangent plane $T_z\qeightLag$ is the graph of the symmetric map
$\qeightHess g(z)\colon \qeightLdual\to L$ (the Hessian of $g$, symmetric by equality of
mixed partials), so the Gauss map of $\qeightLag$ factors through the contractible
graphical locus $\mathcal G$. Hence $\det{}^2\circ(\text{Gauss map})\colon\qeightLag\to
S^1$ factors through a contractible space, its degree on every loop is $0$, and the
Maslov class of $\qeightLag$ vanishes.
\end{proof}

Applying Lemma~\ref{lem:maslov-zero} with $g=f+q_\Sigma$ (smooth for
$f\in\qeightSm(\Sigma)$) shows the smoothed surface $\qeightLag_{\qeightdd f}$ has vanishing Maslov
class on every smoothing chart. Thus, far from being obstructed, the smoothed
Lagrangian is exact (it is the graph of an exact $1$-form $\qeightdd(f+q_\Sigma)$) and
Maslov-trivial; the putative obstruction is identically satisfied.

\paragraph{Where the four-valence enters, and what it is not.} The role of the
$4$-faces hypothesis is to guarantee, via Lemma~\ref{lem:linear}, that the vertex
admits the product-of-two-corners normal form, hence the homeomorphism
$\Sigma\cong\qeightLdual$ of Corollary~\ref{cor:Ldual} and the single-valued
piecewise-quadratic generating function $q_\Sigma=S(Q_1)+S(Q_2)$ with
$S(Q)=\tfrac12 Q|Q|$ (in the symplectic base coordinates $Q_i=(p_i-q_i)/\sqrt2$,
conjugate to $P_i=(p_i+q_i)/\sqrt2$, in which the corner reads $P_i=|Q_i|$). This
is a statement of linear symplectic geometry, not a parity or rotation-number
condition. In particular it is not a Maslov constraint: the obstruction-theoretic
quantities ($\mu$, $tb$, $r$, DGA augmentations) play no role because, as just
shown, the construction never leaves the Maslov-trivial graphical category. The
reconciliation of the two candidate answers is therefore clean: the
YES-construction is correct, and the NO-candidate rests on importing a
filling-category obstruction that does not govern local graphical smoothings.

\section*{Role of the $4$-faces-per-vertex hypothesis}

The hypothesis is used exactly once, in Lemma~\ref{lem:linear}: four Lagrangian
quadrants meeting at a vertex, with edge tangents $(v_1,v_2,u_1,u_2)$, force
$\omega(v_i,u_i)\ne0$ and hence a standard symplectic basis, which is what makes
the symplectic pairing $\Sigma\to\qeightLdual$ a homeomorphism
(Corollary~\ref{cor:Ldual}) and supplies the smooth structure used throughout.
With a different number of faces the local model changes and this normal form is
unavailable; the mechanism is the linear symplectic geometry of the vertex, not a
Maslov-class parity condition.

\section*{Remarks}

\begin{enumerate}\setlength\itemsep{4pt}
\item \textbf{Why smoothing functions, not Legendrian fillings.} The smoothing is
produced directly as graphs over $K$ via the conormal fibration, so no
holomorphic-curve or contact-fillability input is needed; the only analytic
ingredients are Weinstein's neighbourhood theorem and a partition-of-unity
estimate.
\item \textbf{The rescaling subtlety.} $\qeightSm(\Sigma)$ is closed under adding
smooth functions but not under scalar multiplication, so small smoothing
functions must be built by hand (Lemma~\ref{lem:small}); this is the one place a
genuine estimate is required. The smoothing path in Section~6 is therefore the
explicit family $f^\varepsilon\in\qeightSm(K)$ of Lemma~\ref{lem:cont-family}, which
stays inside the admissible class $\qeightSm(K)$ for every $\varepsilon>0$ and whose
$C^1$-norm tends to $0$; we do \emph{not} scale a smoothing function, since that
would leave $\qeightSm(K)$ (the coset $g+C^\infty(K)$ does not contain $0$).
\item \textbf{Matched normals.} Property (iii) of Definition~\ref{def:fibration}
ensures the two smooth structures on $K$ — from Corollary~\ref{cor:Ldual} and from
the edge collars — agree, so the global graph construction is consistent across
edges.
\item \textbf{Generalization.} The argument is local-to-global and uses only
linear symplectic normal forms plus partition-of-unity gluing, so it adapts to
other ambient symplectic manifolds once the vertex normal form is established.
\end{enumerate}

\begin{thebibliography}{9}
\bibitem{Hirsch} M.~W.~Hirsch.
  \emph{Differential Topology}. Graduate Texts in Mathematics, Springer, 1976.
\bibitem{McDuffSalamon} D.~McDuff and D.~Salamon.
  \emph{Introduction to Symplectic Topology}, 3rd ed.,
  Oxford University Press, 2017.
\end{thebibliography}

\endgroup
